% Short report (eight pages) on the main findings of muon/notes_spectral_theory.tex. A longer condensed version is
% muon/paper_spectral_muon.tex; the bibliography file is shared with it.
% Build (from muon/):  pdflatex report_spectral_muon && bibtex report_spectral_muon && pdflatex report_spectral_muon && pdflatex report_spectral_muon
% Figure: .venv/Scripts/python muon/theory_checks/make_paper_figures.py T2s
\documentclass[11pt]{article}
\usepackage[utf8]{inputenc}
\usepackage[T1]{fontenc}
\usepackage{amsmath,amssymb,amsthm,mathtools}
\usepackage{graphicx}
\usepackage{booktabs}
\usepackage{enumitem}
\usepackage[margin=1in]{geometry}
\usepackage[colorlinks=true,linkcolor=blue!55!black,citecolor=blue!55!black,urlcolor=blue!55!black]{hyperref}
\usepackage{xcolor}
\graphicspath{{figures/}}
\setlength{\emergencystretch}{2em}

\newtheorem{theorem}{Theorem}
\newtheorem{proposition}[theorem]{Proposition}
\newtheorem{lemma}[theorem]{Lemma}

\newcommand{\R}{\mathbb{R}}
\newcommand{\E}{\mathbb{E}}
\newcommand{\Prob}{\mathbb{P}}
\newcommand{\norm}[1]{\left\|#1\right\|}
\newcommand{\nop}[1]{\left\|#1\right\|_{\mathrm{op}}}
\newcommand{\nnuc}[1]{\left\|#1\right\|_{*}}
\newcommand{\nF}[1]{\left\|#1\right\|_{F}}
\newcommand{\inner}[2]{\left\langle #1, #2 \right\rangle}
\newcommand{\abs}[1]{\left|#1\right|}
\newcommand{\tr}{\operatorname{tr}}
\newcommand{\diag}{\operatorname{diag}}
\newcommand{\polar}{\operatorname{polar}}
\newcommand{\Var}{\operatorname{Var}}
\newcommand{\MAD}{\operatorname{MAD}}
\newcommand{\Med}{\operatorname{Med}_\ast}
\newcommand{\Lop}{L_{\mathrm{op}}}
\newcommand{\muop}{\mu_{\mathrm{op}}}
\newcommand{\kop}{\kappa_{\mathrm{op}}}
\newcommand{\kF}{\kappa_{F}}
\newcommand{\lmin}{\lambda_{\min}}
\newcommand{\lmax}{\lambda_{\max}}
\newcommand{\ones}{\mathbf{1}}

% compact bibliography
\let\oldthebibliography\thebibliography
\renewcommand{\thebibliography}[1]{\oldthebibliography{#1}\setlength{\itemsep}{0.5pt}\setlength{\parskip}{0pt}}

\title{\vspace{-2.2em}What limits Muon on strongly convex problems?\\ Conditioning barriers and convergence to a nuclear median}
\author{Side study of the project ``Accelerating Variance Reduction''\\ Statistical Learning Theory, UC Berkeley}
\date{\vspace{-0.4em}16 November 2026}

\begin{document}
\maketitle

\begin{abstract}
Muon updates a matrix parameter along the polar factor of its (averaged) gradient: steepest descent for the spectral
norm. Coming from a study of variance-reduced methods, we ask what limits this step on strongly convex problems, in
speed and in precision. \emph{Conditioning.} For ridge-regularized generalized linear models whose loss curvature
vanishes along a sequence of parameters (logistic, softmax), no norm, spectral or other, has a better worst-case
condition number than the Euclidean one. \emph{Noise.} With additive gradient noise, stochastic polar steps are
stochastic mirror descent, in the space of gradients, on a noise-smoothed nuclear norm. For step lengths with
$\sum_t\eta_t=\infty$ and $\sum_t\eta_t^2<\infty$ the gradient converges almost surely to the negative of a
\emph{nuclear median} of the noise: the iterates converge to the minimizer if zero is the only such median, and stay
away from it if zero is not one. For batch size one on least squares the limit is a weighted least-distances
estimator; for skewed noise with independent entries a formal expansion, checked by simulation, gives an offset of
order $1/b$ in the batch size $b$; a finite-sum variance-reduced estimator has the minimizer as its only fixed point.
\emph{Dynamics.} On least squares we add to the known polar dynamics a conservation law in dimension two, an integer
counterexample to the monotonicity of the spectral norm of the error, and a contraction factor at most $1-2/\kop$ for
the second nuclear-scaled step. Published facts are cited as such; the other statements were not found in our
literature checks, and those that were not searched are listed.
\end{abstract}

\section{Introduction}

Muon \cite{jordan2024muon} replaces the gradient $G$ of each weight matrix by an approximation of its polar factor
$\polar(G)=UV^\top$ ($G=U\diag(s)V^\top$): steepest descent for the spectral norm
\cite{carlson2015stochastic,bernstein2024old}. A fast-growing literature analyses this step on models where everything
can be computed: matrix quadratics \cite{shulgin2026quadratic,davis2025spectral}, linear models
\cite{vasudeva2025muon} and noisy gradients \cite{xie2026muon,chen2026when,paquette2026phases,do2026muon}. This
report comes from a different direction: a project on variance-reduced methods for finite sums (SAGA
\cite{defazio2014saga}, SVRG \cite{johnson2013accelerating}, the analysis of \cite{bach2024learning}), which concluded
that worst-case constants hide what happens along trajectories and that no fixed quadratic metric improves the
worst-case condition number of regularized logistic regression. We ask what these conclusions become for Muon, and
what a variance-reduction viewpoint says about its \emph{speed} and its \emph{precision}. A first set of answers went
through an adversarial review of the mathematics and of the literature; a large part turned out to be published in
2025--2026 and is cited as such below. This report states what remains; full proofs and the record of the review are
in a companion note (\texttt{muon/notes\_spectral\_theory.pdf}).

\paragraph{Related work.} Attributions are given next to each statement; we only recall the closest results on
noise. That gradient noise smooths the polar map, and that Muon at low signal-to-noise ratio is a gradient step of
size $\eta\,\E\nnuc Z/(Kd\sigma)$, are in \cite{xie2026muon,chen2026when,paquette2026phases}. That sign methods drive
a median, not the mean, of the stochastic gradient to zero is known \cite{chen2020distributed}; normalized steps on a
location problem estimate the geometric median \cite{cardot2013efficient}; variance reduction inside sign or
steepest-descent steps appears in \cite{chzhen2023signsvrg,li2026implicit}.

\section{Setting and the dual picture}\label{sec:setting}

$W\in\R^{K\times d}$, $r=\min(K,d)$, $\inner AB=\tr(A^\top B)$; $\nF\cdot$, $\nop\cdot$, $\nnuc\cdot$ are the
Frobenius, spectral and nuclear norms. For $M=U\diag(s)V^\top$ (reduced on the nonzero singular values),
$\polar(M)=UV^\top$ and $\abs M=(M^\top M)^{1/2}$. The function $\psi(G)=\nnuc G$ is convex and
$\polar(G)\in\partial\psi(G)$. A \emph{polar step of length $\eta$} is $W^+=W-\eta\,\polar(G)$, $G=\nabla F(W)$; the
\emph{nuclear-scaled step} is $W^+=W-\frac{\nnuc G}{L}\polar(G)$ \cite{carlson2015stochastic,lau2025polargrad}. $F$ is
$C^2$ with minimizer $W_\ast$ and minimum $F_\ast$, and for a norm $N$ we let $\mu_N,L_N$ be the best constants in
$\mu_NN(\Delta)^2\le D^2F(W)[\Delta,\Delta]\le L_NN(\Delta)^2$ and $\kappa_N=L_N/\mu_N$; subscripts ``op'' and $F$
refer to the spectral and Frobenius norms. Two models are used: (LS)
$F(W)=F_\ast+\frac12\tr((W-W_\ast)\Sigma(W-W_\ast)^\top)$, $\Sigma\succ0$, i.e.\ multi-output least squares, with
$E=W-W_\ast$ and $G=E\Sigma$; (GLM) $F(W)=\frac1n\sum_i\ell_i(Wx_i)+\frac\lambda2\nF W^2$ with convex $C^2$ losses,
$c_i=\sup_z\lmax(\nabla^2\ell_i(z))<\infty$, and $f_i(W)=\ell_i(Wx_i)+\frac\lambda2\nF W^2$.

\begin{lemma}[dual picture]\label{lem:dual}
Assume $0<\muop\le\Lop<\infty$ and let $F^\ast$ be the conjugate of $F$, $D_\Phi$ the Bregman divergence of $\Phi$.
(a) A polar step is one step of mirror subgradient descent on $\psi$ in the space of gradients with mirror map
$F^\ast$: $\nabla F^\ast(G^+)=\nabla F^\ast(G)-\eta O$, $O\in\partial\psi(G)$.
(b) $D_{F^\ast}(\nabla F(W'),\nabla F(W))=D_F(W,W')$; in particular $D_{F^\ast}(0,G)=F(W)-F_\ast$.
(c) $2\muop(F(W)-F_\ast)\le\nnuc{\nabla F(W)}^2\le2\Lop(F(W)-F_\ast)$.
\end{lemma}
Item (b) is Lemma~3.2 of \cite{maddison2021dual}, item (c) the dual-norm Polyak--\L{}ojasiewicz inequality and its
converse \cite{balles2020geometry,leplat2025norm}. Methods $W^+=W-\eta\nabla k(\nabla F(W))$ are dual-space
preconditioned gradient descent \cite{maddison2021dual}, read by its authors as a Bregman method in gradient space
for smooth $k$; (a) is the nonsmooth instance $k=\nnuc\cdot$ of this reading. The nuclear-scaled step with $L=\Lop$
is steepest descent for the spectral norm with step $1/\Lop$, with the known rate $1-1/\kop$
\cite{balles2020geometry}.

\section{Conditioning: what no norm can buy}\label{sec:kappa}

\begin{theorem}[no free lunch, for every norm]\label{thm:nfl}
Consider (GLM) and assume that there is a sequence $W_t$ with $\max_i\nop{\nabla^2\ell_i(W_tx_i)}\to0$ (logistic and
softmax losses: $W_t=tW_1$ for generic $W_1$). Let $N$ be any norm on $\R^{K\times d}$ and
$c_N=\max_{\nF\Delta=1}N(\Delta)$. Then $\mu_N(F)=\lambda/c_N^2$ and $L_N(g)\ge L_F(g)/c_N^2$ for every convex $C^2$
function $g$. Consequently
\[
\kappa_N(F)\ \ge\ \kF(F),\qquad \frac{L_N(f_i)}{\mu_N(F)}\ \ge\ \frac{L_F(f_i)}{\mu_F(F)}=\frac{c_i\norm{x_i}^2+\lambda}{\lambda}\quad\text{for every }i .
\]
\end{theorem}
\begin{proof}
For all $W,\Delta$: $D^2F(W)[\Delta,\Delta]\ge\lambda\nF\Delta^2\ge(\lambda/c_N^2)N(\Delta)^2$. At a maximizer
$\Delta_0$ of $N/\nF\cdot$,
$D^2F(W_t)[\Delta_0,\Delta_0]\le\max_i\nop{\nabla^2\ell_i(W_tx_i)}\frac1n\sum_i\norm{\Delta_0x_i}^2+\lambda\nF{\Delta_0}^2
\to\lambda\nF{\Delta_0}^2$, so $\mu_N(F)=\lambda/c_N^2$. For convex $g$,
$D^2g[\Delta,\Delta]/N(\Delta)^2\ge D^2g[\Delta,\Delta]/(c_N^2\nF\Delta^2)$; take suprema.
\end{proof}

The statement covers the spectral norm ($c_N=1$: $\muop=\mu_F=\lambda$ while $\Lop(f_i)=c_i\norm{x_i}^2+\lambda r$),
its preconditioned versions $\nop{S^{1/2}\Delta P^{1/2}}$, the $\ell_\infty$ norm of sign descent, and Euclidean or
Kronecker-factored metrics; for $K=1$ and Euclidean metrics it is the no-free-lunch result of our main project. The
curvature floor of a ridge-regularized model is isotropic, while a norm can only trade smoothness against strong
convexity: any advantage of a non-Euclidean step on such problems is a property of the trajectory, not of the
constants. The hypothesis on the losses cannot be dropped: for ridge least squares the norm
$\nF{\Delta\Sigma^{1/2}}$ has condition number one.

In the spectral geometry there is in addition a dimension barrier, $\kop\ge\min(K,d)$, the matrix form of a folklore
fact on $\ell_\infty$ \cite{sidford2018coordinate}, and finite sums lose a familiar advantage.

\begin{proposition}[no averaging gain]\label{prop:avg}
For (LS) with $K\ge d$, $\Sigma=\hat\Sigma+\lambda I$, $\hat\Sigma=\frac1n\sum_ix_ix_i^\top$, and
$f_i(W)=\frac12\norm{Wx_i-y_i}^2+\frac\lambda2\nF W^2$:
$\Lop(F)=\tr\Sigma=\frac1n\sum_i\Lop(f_i)$, whereas $L_F(F)=\lmax(\Sigma)\le\frac1n\sum_iL_F(f_i)=\tr\hat\Sigma+\lambda$.
\end{proposition}
\begin{proof}
$D^2f_i[\Delta,\Delta]=\norm{\Delta x_i}^2+\lambda\nF\Delta^2\le\norm{x_i}^2+\lambda d$ on the spectral unit ball and
$D^2F[\Delta,\Delta]=\tr(\Delta\Sigma\Delta^\top)\le\tr\Sigma$, with equality in both for every $\Delta$ with
orthonormal columns: an isometry excites all samples at once.
\end{proof}
The value $\Lop(F)=\tr\Sigma$ is known \cite{davis2025spectral}; its reading for finite sums was not found. In the
spectral geometry the smoothness of the average is the average smoothness, so $\kop(F)=\bar L_{\mathrm{op}}/\mu$ with
$\bar L_{\mathrm{op}}=\frac1n\sum_i\Lop(f_i)=\tr\hat\Sigma+\lambda d$ and $\mu=\lmin(\Sigma)=\muop(F)=\mu_F(F)$. This
is the constant
$\bar L_F/\mu$, $\bar L_F=\frac1n\sum_iL_F(f_i)=\tr\hat\Sigma+\lambda$, of variance-reduced methods with importance
sampling \cite{xiao2014proximal}, plus $(d-1)\lambda/\mu$: without ridge, the full-batch nuclear-scaled step has the
condition number of an importance-sampled variance-reduced method while paying $n$ gradients per step. For $K<d$ the
averaging gain reappears.

\section{Gradient noise: convergence to a nuclear median}\label{sec:noise}

Let $\xi$ be a random $K\times d$ matrix and
\[
\psi_\xi(G)=\E\big[\nnuc{G+\xi}-\nnuc\xi\big],\qquad \Phi_\xi(G)=\E\,\polar(G+\xi).
\]
No moment is required. $\psi_\xi$ is convex and $1$-Lipschitz for the nuclear norm, and
$\Phi_\xi(G)\in\partial\psi_\xi(G)$, with $\Phi_\xi(G)=\nabla\psi_\xi(G)$ when $G+\xi$ has full rank almost surely: this
is the randomized-smoothing identity \cite{duchi2012randomized,jiang2026adaptive} for the spectral ball. We call
$\Med(\xi)=\arg\min_C\E[\nnuc{\xi-C}-\nnuc\xi]$ the set of \emph{nuclear medians} of $\xi$ (the median of a law in a
normed space \cite{minsker2015geometric}); $\arg\min\psi_\xi=-\Med(\xi)$ is nonempty, convex and compact, since
$\psi_\xi(G)\ge(2\Prob(\nnuc\xi\le R)-1)\nnuc G-2R$ for any $R$ with $\Prob(\nnuc\xi\le R)>\frac12$.

\begin{theorem}[stochastic polar steps are stochastic mirror descent]\label{thm:smd}
Let $F$ satisfy $0<\muop\le\Lop<\infty$, let $\xi_t$ be i.i.d.\ copies of $\xi$ independent of the past, and
$W_{t+1}=W_t-\eta_t\,\polar(\nabla F(W_t)+\xi_t)$, $G_t=\nabla F(W_t)$. For every $G^\dagger\in\arg\min\psi_\xi$, with
$W^\dagger=\nabla F^\ast(G^\dagger)$:
\begin{enumerate}[label=(\alph*),itemsep=1pt,topsep=2pt]
\item $\E[D_F(W_{t+1},W^\dagger)\mid W_t]\le D_F(W_t,W^\dagger)-\eta_t(\psi_\xi(G_t)-\min\psi_\xi)+\frac{\Lop}2\eta_t^2$;
\item $\E\,\psi_\xi(\bar G_T)-\min\psi_\xi\le\big(D_F(W_0,W^\dagger)+\frac{\Lop}2\sum_{t<T}\eta_t^2\big)/\sum_{t<T}\eta_t$
for the weighted average $\bar G_T=\sum_{t<T}\eta_tG_t/\sum_{t<T}\eta_t$, whatever the lengths;
\item if $\sum_t\eta_t=\infty$ and $\sum_t\eta_t^2<\infty$, $G_t$ converges almost surely to a point of $-\Med(\xi)$;
\item hence, for such lengths, $W_t\to W_\ast$ if $\Med(\xi)=\{0\}$, and if $0\notin\Med(\xi)$ then almost surely
$\lim_tF(W_t)-F_\ast\ge\min_{C\in\Med(\xi)}\nnuc C^2/(2\Lop)>0$.
\end{enumerate}
\end{theorem}
\begin{proof}
With $O_t=\polar(G_t+\xi_t)$, the three-point identity of mirror descent with mirror map $F^\ast$ gives
$D_{F^\ast}(G^\dagger,G_{t+1})=D_{F^\ast}(G^\dagger,G_t)-\eta_t\inner{O_t}{G_t-G^\dagger}+D_{F^\ast}(G_t,G_{t+1})$. By
Lemma~\ref{lem:dual}(b) the first two divergences are $D_F(W_{t+1},W^\dagger)$ and $D_F(W_t,W^\dagger)$, and the last
one is $D_F(W_{t+1},W_t)\le\frac{\Lop}2\nop{W_{t+1}-W_t}^2\le\frac\Lop2\eta_t^2$. Conditionally on $W_t$,
$\E\inner{O_t}{G_t-G^\dagger}=\inner{\Phi_\xi(G_t)}{G_t-G^\dagger}\ge\psi_\xi(G_t)-\psi_\xi(G^\dagger)$. This is (a);
(b) follows by summation and convexity. (c) By the Robbins--Siegmund lemma, almost surely $D_F(W_t,W')$ converges,
simultaneously for a countable dense set of minimizers $\nabla F(W')$, and
$\sum_t\eta_t(\psi_\xi(G_t)-\min\psi_\xi)<\infty$; the iterates are bounded since
$D_F(W_t,W^\dagger)\ge\frac{\muop}2\nop{W_t-W^\dagger}^2$, a cluster point of $G_t$ is a minimizer, and the whole
sequence converges to it. (d) follows from Lemma~\ref{lem:dual}(c).
\end{proof}

The steps being bounded, no variance assumption is needed. The theorem replaces the usual condition $\E\xi=0$ by
$\Med(\xi)=\{0\}$; if $0$ is a nuclear median but not the only one, $W_\ast$ is only one of the possible limits. The
set $-\Med(\xi)$ that contains the limit of the gradient is the same for all such schedules.

\begin{proposition}[medians, zeros of the mean field, size]\label{prop:median}
(a) Every zero of $\Phi_\xi$ belongs to $-\Med(\xi)$; the converse holds at points $G$ where $G+\xi$ has full rank
almost surely, and can fail otherwise.
(b) If $\xi$ and $-\xi$ have the same law, $0\in\Med(\xi)$, and $\Med(\xi)=\{0\}$ if moreover $0$ is in the support
of $\xi$.
(c) If $\xi$ is integrable and $\E\xi=0$, every $C\in\Med(\xi)$ satisfies $\nnuc C\le\E\nnuc{\xi-C}\le\E\nnuc\xi$, so
the limit of Theorem~\ref{thm:smd} has $F(W^\dagger)-F_\ast\le(\E\nnuc\xi)^2/(2\muop)$.
(d) For a SAGA or SVRG estimator, the stochastic gradient at $(W_\ast,\text{exact table})$ is $0$ for every index: a
Dirac mass, whose median is its mean. Such estimators have a fixed point without offset, and it is the only one.
\end{proposition}
\begin{proof}
(a) $\Phi_\xi(G)\in\partial\psi_\xi(G)$, a singleton under the rank condition. (b) $\psi_\xi$ is even and convex, and
$\psi_\xi(G)-\psi_\xi(0)=\frac12\E[\nnuc{\xi+G}+\nnuc{\xi-G}-2\nnuc\xi]$ has a nonnegative integrand, positive on
$\{\nnuc\xi<\nnuc G\}$. (c) $\nnuc C=\nnuc{\E(C-\xi)}\le\E\nnuc{\xi-C}\le\E\nnuc\xi$ by Jensen's inequality and the
minimality of $C$; conclude with Lemma~\ref{lem:dual}(c). (d) Direct substitution.
\end{proof}
Item (b) is the matrix form of the symmetric-noise condition under which sign methods are known to converge
\cite{bernstein2019majority}. Item (d) identifies a fixed point and says nothing about its attraction. When the law
of the stochastic gradient depends on $W$ (minibatches of a finite sum), Theorem~\ref{thm:smd} does not apply as such,
but the same mechanism can be followed by hand in the simplest case.

\begin{proposition}[batch one on least squares]\label{prop:lad}
Let $g_j(W)=(Wx_j-y_j)x_j^\top$ (no ridge term), $x_j\ne0$, $j$ uniform in $\{1,\dots,n\}$, and
$\Lambda(W)=\frac1n\sum_i\norm{Wx_i-y_i}/\norm{x_i}$. Then $\polar(g_j(W))\in\partial_W(\norm{Wx_j-y_j}/\norm{x_j})$,
with Frobenius norm at most one: the batch-one polar step is the stochastic subgradient method on $\Lambda$. With
$\sum\eta_t=\infty$, $\sum\eta_t^2<\infty$ it converges almost surely to a minimizer of the weighted least-distances
objective $\Lambda$ (least absolute deviations for $K=1$), which in general is not the least-squares solution.
\end{proposition}
\begin{proof}
For $r=Wx-y\ne0$, $\polar(rx^\top)=rx^\top/(\norm r\norm x)=\nabla_W\norm{Wx-y}/\norm x$
\cite[Lemma~H.1]{li2026implicit}; for $r=0$, $\polar(0)=0\in\partial_W(\norm{Wx-y}/\norm x)$. Hence
$\E_j\polar(g_j(W))\in\partial\Lambda(W)$ for every $W$, and the rest is the standard analysis of the stochastic
subgradient method in the Frobenius geometry.
\end{proof}
For coordinatewise sign descent with batch size one, non-convergence to the minimizer for every step-size sequence is
proved in \cite{karimireddy2019error} (Theorem~I) when all data points share a sign pattern, with a least-squares
example. For $K\ge2$ the statement was not found.

\begin{proposition}[size of the offset for minibatches; formal]\label{prop:offset}
Let $\xi_b$ be the mean of $b$ i.i.d.\ centred matrices with i.i.d.\ entries of variance $\sigma^2$ and third
cumulant $\kappa_3$, with a non-lattice law. A one-term Edgeworth expansion gives, for the minimizer $G_b^\dagger$ of
$\psi_{\xi_b}$,
\begin{equation}\label{eq:offset}
G_b^\dagger=\frac{\kappa_3}{6\sigma^2b}\cdot\frac{a_{K,d}}{c_{K,d}}\;\ones\ones^\top+o\Big(\frac1b\Big),\qquad
a_{K,d}=-\E\big[\polar(Z)_{11}(Z_{11}^3-3Z_{11})\big],\quad c_{K,d}=\frac{\E\nnuc Z}{Kd},
\end{equation}
with $Z$ a standard Gaussian $K\times d$ matrix: in every entry, the scalar offset $\kappa_3/(6\sigma^2b)$ of the
median of a mean (Cornish--Fisher \cite{cornish1938moments}; $a_{1,1}=c_{1,1}=\sqrt{2/\pi}$) times $a_{K,d}/c_{K,d}$.
The objective gap at the limit is of order $1/b^2$.
\end{proposition}
\begin{proof}[Derivation]
$\polar$ is homogeneous of degree $0$, so $\Phi_{\xi_b}(0)=\E\polar(\sqrt b\,\xi_b)$; the density of $\sqrt b\,\xi_b$
is the Gaussian one times $1+b^{-1/2}\frac{\kappa_3}{6\sigma^3}\sum_{ij}\mathrm{He}_3(z_{ij}/\sigma)+O(1/b)$, and the
even terms integrate to zero against the odd map $\polar$:
$\Phi_{\xi_b}(0)=-b^{-1/2}\frac{\kappa_3}{6\sigma^3}a_{K,d}\ones\ones^\top+O(b^{-3/2})$. The derivative of
$\Phi_{\xi_b}$ at $0$ is $\sqrt b\,(c_{K,d}/\sigma)\mathrm{Id}$ to leading order \cite{xie2026muon,chen2026when}.
Solve $\Phi_{\xi_b}(G)=0$ to first order.
\end{proof}

The expansion is formal (an Edgeworth expansion against a discontinuous map), but its predictions are parameter-free
and are confirmed numerically in Section~\ref{sec:num}. Two readings. First, \emph{what matters is the skewness of
the gradient noise}: in a regression with symmetric inputs independent of the residuals, the per-sample gradient
$\varepsilon x^\top$ is symmetric in law whatever the skewness of $\varepsilon$, and the population offset is zero;
an intercept is enough to create one. Second, $a_{K,d}/c_{K,d}$ is of order $1/d$: $a_{1,d}/c_{1,d}=3/(d+2)$
exactly, a quadrature gives $d\,a_{K,d}/c_{K,d}=1.70$ for $(K,d)=(2,3)$, and Monte Carlo (relative standard errors
below $3\%$) gives $1.37$, $1.71$, $2.00$, $2.17$, $2.29$ for $K=d=2,4,8,16,32$ and $2.59$, $2.80$ for
$(K,d)=(4,16)$, $(16,64)$. Per entry, the
offset of the polar step is thus smaller than that of sign descent, $\kappa_3/(6\sigma^2b)$ in every coordinate, by a
factor of order $d$: orthogonalization averages the skewness. In dimension one and for the Newton--Schulz map with
momentum, non-convergence for bounded data with a nonzero third \emph{central} moment, for almost all batch sizes, is
proved in \cite{do2026muon} (Theorem~1.3).

\section{Least-squares dynamics: complements}\label{sec:ls}

For (LS) with $K\ge d$ the polar dynamics are largely understood \cite{shulgin2026quadratic,vasudeva2025muon}: at
full-rank iterates the nuclear-scaled step with $L=\tr\Sigma=\Lop$ is an exact line search; in the \emph{commuting
case} ($E_0^\top E_0$ commutes with $\Sigma$, so that $E=\sum_ka_ku_kv_k^\top$ with $(v_k)$ an eigenbasis of
$\Sigma$), while the amplitudes $a_k$ are positive, this step maps them to $\abs{a-\E_p[a]}$, the mean being taken
under the weights $p_k=\lambda_k(\Sigma)/\tr\Sigma$, with factor $\Var_p(a)/\E_p[a^2]$ on $F-F_\ast$. The
\emph{polar flow} is the solution of
$\dot W\in-\partial\psi(\nabla F(W))$, unique for (LS) (in the variable $G=E\Sigma$ it is the gradient flow of the
convex function $\psi$ for the inner product $\tr(A\Sigma^{-1}B^\top)$ \cite{brezis1973operateurs}); while $E$ has
rank $d$ it reads $\dot E=-\polar(E\Sigma)$.

\begin{proposition}[complements on (LS), $K\ge d$]\label{prop:ls}
Let $E$ have rank $d$, with singular values $s_1\ge\dots\ge s_d$.
\begin{enumerate}[label=(\alph*),itemsep=1pt,topsep=2pt]
\item \emph{Rotation.} $\polar(E\Sigma)=\polar(E)\,Q$ with $Q=\polar(\abs E\Sigma)\in SO(d)$, and $Q=I$ iff $\abs E$
commutes with $\Sigma$. Along the flow, each simple singular value $s_i$ of $E$, with right singular vector $w_i$,
has $\dot s_i=-w_i^\top Qw_i$.
\item \emph{Conservation law.} If $K=d=2$, every polar step $E^+=E-\eta\,\polar(E\Sigma)$, of any length, conserves
$s_1(E)-s_2(E)$ as long as $\det E$ keeps its sign.
\item \emph{Non-monotonicity.} If $d\ge3$, $\nop E$ can increase along the flow while $F$ and $\nF E$ decrease: for
$E_0=\diag(300,200,1)$ and $\Sigma=\left(\begin{smallmatrix}2&4&12\\4&9&20\\12&20&89\end{smallmatrix}\right)$,
$Q_{11}=-0.0143$, and $\nop{E(t)}$ rises from $300$ to $303.27$ before decreasing. This needs an ill-conditioned
error: for $d=3$, $\frac d{dt}\nop E<0$ for every $\Sigma$ as soon as $s_1^2<(2+\sqrt3)^3s_2s_3\approx52\,s_2s_3$.
\item \emph{Contraction after the first step.} Let $p_{(1)}\le p_{(2)}$ be the two smallest weights
$\lambda_k(\Sigma)/\tr\Sigma$. From every full-rank start, commuting or not, the second nuclear-scaled step
($L=\tr\Sigma$, $W_1\ne W_\ast$) satisfies
\[
\frac{F(W_2)-F_\ast}{F(W_1)-F_\ast}\ \le\ 1-\frac{4p_{(1)}p_{(2)}}{p_{(1)}+p_{(2)}}\ \le\ 1-\frac2{\kop}.
\]
In the commuting case, as long as all amplitudes are positive, the same bound holds at every step $t\ge1$, where the
factor equals $1-\MAD_p(a(t-1))^2/\Var_p(a(t-1))$, $\MAD_p$ being the mean absolute deviation around the mean.
\end{enumerate}
\end{proposition}
\begin{proof}[Sketch]
(a) $\polar(E)Q=E\abs E^{-1}\abs E\Sigma\abs G^{-1}=G\abs G^{-1}$ and $\det(\abs E\Sigma)>0$; $\dot s_i=u_i^\top\dot Ew_i$.
(b) The singular values of a $2\times2$ matrix are $m_c+m_a$ and $\abs{m_c-m_a}$, where $\sqrt2\,m_c$ and
$\sqrt2\,m_a$ are the Frobenius norms of its components in the planes spanned by the rotations and by the reflections.
If $\det E>0$ the polar factor of $E\Sigma$ is a rotation, so the step leaves the second component of $E$ unchanged
and $s_1-s_2=2m_a$; symmetrically if $\det E<0$. (c) Direct computation, the sign of $Q_{11}$ being certified in exact
rational arithmetic; for the last claim see the companion note. (d) After one line-search step
$G_1=\polar(G_0)M_1$ with $M_1$ symmetric of trace zero \cite{shulgin2026quadratic}. With $M_1=M_+-M_-$ and
$S=\tr M_+=\tr M_-$: $\nnuc{G_1}=2S$ and
$2(F(W_1)-F_\ast)=\tr(\Sigma^{-1}M_+^2)+\tr(\Sigma^{-1}M_-^2)\le S^2(1/\lambda_d+1/\lambda_{d-1})$ by convexity and
Ky Fan's inequality; conclude with $F(W_2)\le F(W_1)-\nnuc{G_1}^2/(2\tr\Sigma)$. The commuting case is similar, with
$a(t)=\abs X$ and $X$ centred under $p$.
\end{proof}
The ``if'' part of $Q=I$ is in \cite{shulgin2026quadratic}, and (b) explains the $45^\circ$ trajectories of singular
values observed there for $2\times2$ problems. For $d=2$ the bound of (d) is the exact factor
$((\kappa-1)/(\kappa+1))^2$, $\kappa=\lambda_1/\lambda_2$, proved in \cite{shulgin2026quadratic} for every $t\ge1$ and
every full-rank start; for $d\ge3$ it was not found, and at $t=0$ the factor can be arbitrarily close to the certified
$1-1/\kop$ \cite{shulgin2026quadratic}.

\section{Numerical illustrations}\label{sec:num}

\begin{figure}[t]
\centering
\includegraphics[width=0.84\textwidth]{figT2s_noise}
\caption{Stochastic polar steps on a finite-sum least-squares problem ($n=2000$, $K=2$, $d=3$, intercept, skewed
residuals). (a) Batch one, same decreasing lengths: with the minibatch gradient the averaged iterate converges to the
least-distances solution $W_{\mathrm{LAD}}$ of Proposition~\ref{prop:lad}; with the SAGA estimator it converges to the
least-squares solution $W_{\mathrm{LS}}$. (b) Batch size times the nuclear norm of the gradient at the limit point:
zeros of the mean field (line), stochastic runs with decreasing lengths (dots), SAGA estimator with the nuclear-scaled
length (squares, at rounding level).}\label{fig:noise}
\end{figure}

\paragraph{Offset.} For $2\times3$ matrices with entries $\mathrm{Exp}(1)-1$ ($\sigma=1$, $\kappa_3=2$),
\eqref{eq:offset} predicts $\sqrt b\,\|\Phi_{\xi_b}(0)\|_F\to0.234$ and $b\,\|G_b^\dagger\|_*\to0.464$. Computing the
zero of the mean field directly, with exactly sampled batch means and $1.5\cdot10^6$ samples, gives $0.230$--$0.237$
and $0.450$--$0.498$ for $b=2,\dots,256$. The fixed point is proportional to the all-ones matrix (cosine above
$0.999$ for $b\le64$), and with symmetrized noise the mean field at $0$ vanishes up to the Monte-Carlo error
($10^{-3}$).

\paragraph{Finite sum.} The dichotomy of Theorem~\ref{thm:smd} is visible on a finite sum
(Figure~\ref{fig:noise}): a least-squares problem with an intercept and skewed residuals, for which
$\nF{W_{\mathrm{LAD}}-W_{\mathrm{LS}}}=0.170$. Batch-one polar steps with decreasing lengths end (iterates averaged
over the last half of the run and over $32$ runs) at distance $2.6\cdot10^{-3}$ from the least-distances solution and
$0.169$ from the least-squares solution, while the same lengths with the SAGA estimator end at $2.1\cdot10^{-4}$ from
the least-squares solution. For $b=1,2,4,8,16,32$ the zero of the mean field has $b\,\nnuc{\nabla F}$ equal to
$0.168$, $0.240$, $0.276$, $0.285$, $0.303$, $0.308$, and the stochastic runs (same averaging, $64$ runs) end within
$5\%$ of these values; with the SAGA estimator and the nuclear-scaled length the gradient norm ends below
$2.4\cdot10^{-15}$ for every $b$. In this design the per-sample gradients have rank one and the entries of the noise
are not independent, so \eqref{eq:offset} does not apply as such: what is observed is the $1/b$ regime, a plateau of
$b$ times the offset.

\section{Discussion}\label{sec:discussion}

\paragraph{What this says about Muon.} On the side of \emph{speed}, the spectral geometry cannot improve a worst-case
constant on ridge-regularized classification models (Theorem~\ref{thm:nfl}), and its condition number is at least
$\min(K,d)$; what polar steps buy is a property of trajectories. On the side of \emph{precision}, besides the floor
of constant step lengths and the fluctuations of any constant-step stochastic method, a polar step with stochastic
gradients faces the median offset of Theorem~\ref{thm:smd}, whose set of possible limits is the same for all
schedules with $\sum_t\eta_t=\infty$, $\sum_t\eta_t^2<\infty$. When the gradient noise is symmetric, $W_\ast$ is a
possible limit, and the only one if $0$ is in the support of the noise (Proposition~\ref{prop:median}(b)). Otherwise
the offset is formally of order $1/b$ when the noise has a nonzero third cumulant (Proposition~\ref{prop:offset}),
and smaller per entry than for sign descent by a factor of order $d$. A finite-sum variance-reduced estimator has
the minimizer as its only fixed point: this is the precise sense in which variance reduction can improve the
precision of Muon-type steps, not by reducing a variance term in a rate but by changing the fixed point (the
convergence to it is observed in Figure~\ref{fig:noise}, not proved).

\paragraph{Limitations and status.} All statements are for strongly convex objectives. Theorem~\ref{thm:smd} assumes
additive noise and no momentum; for minibatches of a finite sum the law of the noise depends on the iterate, and we
only treat batch one on least squares. Proposition~\ref{prop:offset} is a formal expansion, supported by simulations.
No convergence of a variance-reduced polar method is proved: Proposition~\ref{prop:median}(d) only identifies its
fixed point. The scaling of $a_{K,d}/c_{K,d}$ in $1/d$ is numerical (exact for $K=1$), and it concerns noise with
independent entries in the regime $b\to\infty$ at fixed size. ``Not found'' refers to automated literature checks of
12 and 16 November 2026, which opened about sixty papers; several of the closest ones are a few weeks old. Three
statements were added after these checks and were not searched separately: Theorem~\ref{thm:nfl} for an arbitrary
norm (the checks covered the spectral norm and its preconditioned versions), the condition for $d=3$ in
Proposition~\ref{prop:ls}(c), and the bound of Proposition~\ref{prop:ls}(d) for non-commuting starts.

\paragraph{Reproducibility and use of AI tools.} Every number comes from the scripts of the folder
\texttt{muon/theory\_checks}. A large language model (Claude, Anthropic) was used, under the supervision of the
authors, to derive and check the statements, to write the scripts and to run the adversarial review; statement
numbers quoted from preprints are those read during that review.

{\footnotesize
\bibliographystyle{abbrv}
\bibliography{paper_spectral_muon}
}

\end{document}

